\documentclass[12pt,a4paper]{article}

% --- Преамбула ---
\usepackage[utf8]{inputenc}
\usepackage[T1]{fontenc}
\usepackage[russian]{babel}
\usepackage{amsmath,amssymb}
\usepackage{hyperref}
\usepackage[margin=2.5cm]{geometry}
\usepackage{enumitem}
\usepackage{setspace}
\usepackage{booktabs}
\onehalfspacing

\title{\textbf{Третье дополнение к Субквантовой теории}\\
\large Принцип мгновенной сообщаемости через общее начало}
\author{}
\date{Версия 1.1}

\begin{document}
\maketitle

\section{Введение}
В предыдущих дополнениях были введены понятия ортогональности, константы масштаба и принципа масштабной проницаемости. В данном дополнении формализуется принцип мгновенной сообщаемости --- механизм, позволяющий двум системам с идентичной топологией обмениваться информацией без задержки, независимо от расстояния в эмерджентном пространстве-времени.

Этот принцип является прямым следствием структуры собственных мод и существования общего корня всех мод в Мультиверсе.

\section{Теоретические основы}

\subsection{Общее начало}
Пусть две частицы $A$ и $B$ имеют собственные моды $\psi_A(k)$ и $\psi_B(k)$, определённые на дискретной оси $k = 0, 1, 2, \dots, d_{AA}$ (или $d_{BB}$).

\begin{quote}
\textbf{Определение 1 (Общее начало).} Моды $\psi_A$ и $\psi_B$ имеют \textbf{общее начало} на масштабе $K$, если для всех $k \in [0, K]$ выполняется:
\begin{equation}
    \psi_A(k) = \psi_B(k)
\end{equation}
с точностью до глобальной фазы. Это означает, что формы мод идентичны на начальном участке длины $K$.
\end{quote}

Частицы, имеющие общее начало на масштабе $K$, «выросли» из одной топологической конфигурации. Их causal history совпадает на этом масштабе.

\subsection{Сонаправленность и встречность}
В общем случае моды двух частиц могут находиться в двух типах отношений:
\begin{itemize}
    \item \textbf{Встречные моды.} Моды направлены друг на друга: $\psi_A(k) = -\psi_B(d - k)$ для некоторого расстояния $d$. Это стандартное взаимодействие через посредников (калибровка $F_{A \to B}(0) = -F_{B \to A}(0)$).
    \item \textbf{Сонаправленные моды.} Моды направлены в одну сторону: $\psi_A(k) = \psi_B(k)$ для $k \in [0, K]$. Это общее начало. Взаимодействие через общий корень, без посредников. Именно этот случай ответственен за мгновенную сообщаемость.
\end{itemize}

\subsection{Общий корень}
Все собственные моды всех частиц сети $\Omega$ сходятся в единой точке при $k \to \infty$, называемой общим корнем. Значение поперечного спектра в общем корне есть универсальная константа Мультиверса $F_0$.

Общий корень гарантирует, что любые две моды «встречаются» в бесконечности. Но корреляция через общий корень максимальна именно для тех мод, которые имеют общее начало --- то есть совпадают на конечном интервале.

\section{Аксиома сообщаемости}

\begin{quote}
\textbf{Аксиома 1 (Мгновенная сообщаемость).} Если две частицы $A$ и $B$ имеют общее начало на масштабе $K$, то любое изменение собственной моды частицы $A$ на интервале $k \in [0, K]$ мгновенно вызывает соответствующее изменение собственной моды частицы $B$ на том же интервале.
\end{quote}

Изменение является обратным (с обратным знаком) по отношению к изменению, вызвавшему его. Это обеспечивает сохранение глобальной causal history сети.

Для наблюдателя, не знающего о существовании общего начала, это изменение выглядит как случайная корреляция.

\section{Коэффициент сообщаемости и проблема множественных копий}

\subsection{Запрет на точные копии}
В силу бесконечного пространства вариантов (бесконечной вложенности) точные копии частиц запрещены. Causal history двух частиц одного типа никогда не бывает полностью идентичной. Всегда существуют микроскопические отличия, накопленные за их индивидуальную историю.

Поэтому общее начало всегда существует лишь на конечном масштабе $K$. Для двух электронов в одинаковом состоянии $K$ может быть очень велико, но не бесконечно.

\subsection{Коэффициент сообщаемости}

\begin{quote}
\textbf{Определение 2 (Коэффициент сообщаемости).} Пусть $K_{\max}$ --- максимальный масштаб, на котором моды $\psi_A$ и $\psi_B$ совпадают. Тогда коэффициент сообщаемости:
\begin{equation}
    \chi_{AB} = \frac{K_{\max}}{\min(d_{AA}, d_{BB})}
\end{equation}
\end{quote}

\begin{itemize}
    \item $\chi \to 1$: практически полная сообщаемость (идентичные макрообъекты --- антенны).
    \item $\chi = 0$: сообщаемость отсутствует (разная топология нулей).
    \item $0 < \chi < 1$: частичная сообщаемость.
\end{itemize}

\subsection{Размазывание сигнала}
Если существует $N$ частиц, имеющих общее начало на масштабе $K$ с частицей $A$, то изменение causal history частицы $A$ распределяется между всеми $N$ частицами. Амплитуда изменения, приходящаяся на каждую частицу $B_i$, обратно пропорциональна $N$:
\begin{equation}
    \Delta\psi_{B_i} \approx \frac{\Delta\psi_A}{N}
\end{equation}

Для одиночного электрона ($N \sim 10^{80}$ во Вселенной) изменение, передаваемое на другой конкретный электрон, чудовищно мало и практически ненаблюдаемо.

Для макрообъекта (антенны), существующего в двух экземплярах ($N = 2$), изменение передаётся полностью. Именно поэтому для практической связи необходимо, чтобы антенны были не только идентичны, но и \textbf{уникальны} --- чтобы никакая третья система не имела такого же общего начала и не «размазывала» сигнал.

\section{Устойчивость сообщаемости для макрообъектов}
В отличие от одиночных квантовых частиц, макрообъекты (антенны) обладают свойством \textbf{неразрушающего наблюдения}.

\begin{itemize}
    \item \textbf{Квантовая частица (электрон):} Измерение меняет её моду на масштабе $K$. Общее начало с другой частицей разрушается. Сообщаемость теряется после однократного измерения (коллапс волновой функции).
    \item \textbf{Макрообъект (антенна):} Измерение (считывание сигнала) затрагивает лишь малую часть causal history антенны. Масштаб $K$ общего начала настолько велик, что локальное возмущение не разрушает идентичность мод на всём масштабе. Поэтому сообщаемость сохраняется, и канал связи может использоваться многократно.
\end{itemize}

Это делает топологические антенны принципиально отличными от запутанных частиц: они позволяют передавать данные более одного раза без необходимости перезапутывания.

\section{Связь с известными явлениями}

\subsection{Квантовая запутанность}
Две запутанные частицы имеют общее начало на масштабе $K$, равном их совместной causal history. Измерение одной меняет другую мгновенно. Но после измерения общее начало разрушается --- повторная передача невозможна.

\subsection{Топологические антенны}
Два идентичных кристалла, вырезанные из одной пластины, имеют $\chi \approx 1$ и уникальный рисунок ($N = 2$). Изменение causal history одного кристалла (модуляция сигнала) мгновенно отражается на другом. Наблюдение не разрушает связь.

\subsection{Универсальность законов физики}
Все электроны имеют общее начало на масштабе $K_{\text{common}}$ (общая топология --- $2$ нуля). Поэтому их свойства (масса, заряд) одинаковы. Изменение causal history одного электрона распределяется на все электроны Вселенной, но амплитуда изменения для каждого ничтожно мала. Именно поэтому законы физики стабильны: чтобы изменить массу электрона, нужно изменить causal history значительной доли всех электронов одновременно.

\section{Предсказания}
\begin{enumerate}[label=\arabic*.]
    \item \textbf{Топологическая связь.} Две антенны с идентичной топологией и уникальным рисунком должны демонстрировать мгновенную корреляцию сигналов, не зависящую от расстояния.
    \item \textbf{Зависимость от числа копий.} Если изготовить $N$ идентичных антенн, амплитуда сигнала между любой парой должна падать как $1/N$.
    \item \textbf{Неразрушающее считывание.} Многократная передача данных через одну и ту же пару антенн должна быть возможна без деградации сигнала.
    \item \textbf{Квантовый предел.} Для одиночных частиц (электронов) мгновенная сообщаемость должна существовать, но быть практически ненаблюдаемой из-за огромного $N$.
\end{enumerate}

\section{Заключение}
Принцип мгновенной сообщаемости через общее начало является прямым следствием структуры собственных мод и существования общего корня. Он объясняет квантовую запутанность, универсальность законов физики и работу топологических антенн. Ключевыми факторами для практического использования являются уникальность системы и малое число копий, чтобы избежать размазывания сигнала. Дополнение готово для включения в Субквантовую теорию.

\end{document}